% Propietario conservado: /Users/ruben/Documents/New project/output/EXTREMA_MEDIA_RAZON_RECIPROCIDAD_ES_EN_20260918/ARTICULO/ES/source/nuclear/08.tex
\chapter{Conservation of action forms and oriented area in the tritic regimes}
\label{x_reciprocidad:nuclear:8}
\section{Balance between two action forms}

The angular outputs \(A,C^*\), in a common unit, determine the form \(\eta=\operatorname{artanh}(C^*/A)\) on the domain \(A\ne0\), \(|C^*/A|<1\). Given two successive forms \(\eta_n,\eta_{n+1}\), define

\[
M_n=\operatorname{diag}(e^{\eta_n/2},e^{-\eta_n/2}),\qquad
G_n=M_n^{-\mathsf T}M_n^{-1},\qquad
I_n(z)=\tfrac12z^\mathsf TG_nz.
\]

Here \(z=(Q,P)^\mathsf T\), with both coordinates having the dimension of the square root of action. The scale \(\hbar\), when it fixes the contour \(I_n=\hbar\), is the previously constructed action output; it does not serve as a fitting target.

Let \(R_n\) be the canonical rotation of angular advance \(\theta_n\). The transport without angular advance and the transport with that advance are, respectively,

\begin{equation}
B_n=M_{n+1}M_n^{-1},\qquad
C_n=M_{n+1}R_nM_n^{-1}.
\label{x_reciprocidad:nuc08:eq:1}
\end{equation}

Both go from chart \(n\) to chart \(n+1\). Define the compressed reading and its normalized complementary register:

\begin{equation}
T_n=\frac{8B_n+C_n}{9},\qquad
D_n=\frac{\sqrt8}{9}(C_n-B_n).
\label{x_reciprocidad:nuc08:eq:2}
\end{equation}

The operator \(D_n\) is an isometrically normalized realization of the unique pattern of differences among the nine phases. The nine-component register is \((8z,-z,\ldots,-z)/27\), with \(z=(C_n-B_n)v\); the sum of its inner products gives the same factor \(8/81\). The nine-phase register is retained when the individual labels enter into other readers.

\begin{theorem}
Action is conserved between reading and memory:

\begin{equation}
\boxed{G_n=T_n^\mathsf TG_{n+1}T_n+D_n^\mathsf TG_{n+1}D_n.}
\label{x_reciprocidad:nuc08:eq:3}
\end{equation}
\end{theorem}

\begin{proof}
Normalization of the charts gives
\(B_n^\mathsf TG_{n+1}B_n=G_n=C_n^\mathsf TG_{n+1}C_n\).
On expanding \eqref{x_reciprocidad:nuc08:eq:2}, the cross contributions have coefficients \(+8\) and \(-8\). The remaining diagonal coefficients are \(72\) for \(B_n\) and \(9\) for \(C_n\), divided by \(81\). The sum is \(G_n\).
\end{proof}

With \(z_{n+1}=T_nz_n\), \(m_n=D_nz_n\), the scalar form is

\begin{equation}
\boxed{I_n(z_n)=I_{n+1}(z_{n+1})+I_{n+1}(m_n).}
\label{x_reciprocidad:nuc08:eq:4}
\end{equation}

The telescoping sum, retaining each memory in the chart where it was emitted, gives

\begin{equation}
I_0(z_0)=I_N(z_N)+\sum_{n=0}^{N-1}I_{n+1}(m_n).
\label{x_reciprocidad:nuc08:eq:5}
\end{equation}

The transport \(B_n\) is an effective part of this composition: replacing it by the identity when the forms differ changes the domains and may invalidate \eqref{x_reciprocidad:nuc08:eq:3}.

\section{Exact accumulation law with variable phase}

In normalized coordinates \(y_n=M_n^{-1}z_n\), the reading is \(y_{n+1}=(8I+R_n)y_n/9\). For a planar rotation,

\begin{equation}
\left(\frac{8I+R_n}{9}\right)^\mathsf T
\left(\frac{8I+R_n}{9}\right)=r_n I,\qquad
r_n=1-\frac{32}{81}\sin^2(\theta_n/2).
\label{x_reciprocidad:nuc08:eq:6}
\end{equation}

Therefore,

\begin{equation}
I_N(z_N)=I_0(z_0)\prod_{n=0}^{N-1}r_n.
\label{x_reciprocidad:nuc08:eq:7}
\end{equation}

\begin{theorem}
For \(I_0(z_0)>0\), the terminal component tends to zero exactly when

\begin{equation}
\sum_{n\ge0}\sin^2(\theta_n/2)=+\infty.
\label{x_reciprocidad:nuc08:eq:8}
\end{equation}
\end{theorem}

\begin{proof}
Let \(c=32/81\) and \(u_n=\sin^2(\theta_n/2)\in[0,1]\). We have \(cu_n\leq-\log(1-cu_n)\leq cu_n/(1-c)\), by integrating \(1/(1-s)\) between zero and \(cu_n\). The sum of these logarithms diverges exactly when \(\sum u_n\) diverges. Applying the logarithm to the product in \eqref{x_reciprocidad:nuc08:eq:7} gives the equivalence.
\end{proof}

Equation \eqref{x_reciprocidad:nuc08:eq:5} is preserved in both cases. A sequence of angular advances may leave a positive terminal component or transfer all action to the register, depending on the sum in \eqref{x_reciprocidad:nuc08:eq:8}. Change of form alone preserves action; angular dispersion determines its allocation in this realization. The phases \(\theta_n\) must come from the specific reader of the trajectory: the proof does not retrospectively select a sequence to impose a result.

\section{Extended Euler and common oriented area}

The generators of the three quadratic regimes are

\[
J_{+1}=(2C_{A_2}+I)/\sqrt3,\qquad
J_0=N,\qquad
J_{-1}=(2F_{\rm av}^2-3I)/\sqrt5.
\]

In their two-dimensional charts they have trace zero and satisfy
\(J_\tau^2=-\tau I\). The Coxeter element comes from the APP antecedent, the nilpotent from threshold transport and \(F_{\rm av}\) from the descent of TPK modes. The exponential functions constitute the subsequent analytic realization:

\[
E_\tau(t)=\exp(tJ_\tau)=c_\tau(t)I+s_\tau(t)J_\tau,
\]

where \((c_\tau,s_\tau)\) are \((\cos,\sin)\), \((1,t)\), and \((\cosh,\sinh)\), respectively. Zero trace gives determinant one. With

\[
\Omega=\begin{pmatrix}0&1\\-1&0\end{pmatrix},
\]

we obtain \(E_\tau^\mathsf T\Omega E_\tau=\Omega\). In variable charts we take
\(C_n=M_{n+1}E_{\tau_n}(t_n)M_n^{-1}\) and the same \(B_n\) as in \eqref{x_reciprocidad:nuc08:eq:1}.

\begin{theorem}
The three regimes satisfy

\begin{equation}
\boxed{\Omega=T_n^\mathsf T\Omega T_n+D_n^\mathsf T\Omega D_n,}
\label{x_reciprocidad:nuc08:eq:9}
\end{equation}

with \(T_n,D_n\) from \eqref{x_reciprocidad:nuc08:eq:2}.
\end{theorem}

\begin{proof}
In dimension two, \(L^\mathsf T\Omega L=(\det L)\Omega\). The determinants of \(M_n\), \(B_n\), and \(C_n\) are one. The expansion used in \eqref{x_reciprocidad:nuc08:eq:3}, applied to the alternating form, proves \eqref{x_reciprocidad:nuc08:eq:9}.
\end{proof}

In a fixed chart, the two determinants are

\begin{equation}
\det T_\tau=\frac{65+16c_\tau(t)}{81},\qquad
\det D_\tau=\frac{16(1-c_\tau(t))}{81}.
\label{x_reciprocidad:nuc08:eq:10}
\end{equation}

This follows from \(\det(aI+C)=a^2+a\operatorname{tr}C+\det C\) and \(\operatorname{tr}E_\tau=2c_\tau\). Their sum is one. For the representatives \(C_{A_2}\), \(I+N\), \(F_{\rm av}^2\), the pairs are \((19/27,8/27)\), \((1,0)\), \((89/81,-8/81)\).

The parabolic regime may produce a nonzero rank-one memory with zero area. The hyperbolic regime preserves oriented area through a complement of opposite sign. The elliptic regime additionally permits a positive action balance between the forms of the two fibres. The common alternating form and the metric forms of each regime retain their respective types.

\subsection{Integrated action and orientation}

Let \(\gamma\) be a closed piecewise \(C^1\) curve in the canonical plane. For any real matrix \(L\) of order two,

\[
\oint_{L\gamma}P\,dQ=(\det L)\oint_\gamma P\,dQ.
\]

Indeed, writing \(Q'=aQ+bP\), \(P'=cQ+dP\), the integrals of \(Q\,dQ\) and \(P\,dP\) are zero, and \(\oint Q\,dP=-\oint P\,dQ\). The coefficient \(ad-bc\) remains. Applying \eqref{x_reciprocidad:nuc08:eq:10},

\begin{equation}
\boxed{\oint_\gamma P\,dQ
=\oint_{T_n\gamma}P\,dQ+\oint_{D_n\gamma}P\,dQ.}
\label{x_reciprocidad:nuc08:eq:11}
\end{equation}

The signs are those of the transported curves. The identity is preserved for the three regimes and for changes of form with determinant one.

\section{The common factor of the exponential reader}

The radional reader in \cref{x_reciprocidad:iv:ellipse:section,x_reciprocidad:euler-trit-evaluaciones} also contains \(e^{-x}\exp(-yJ_\tau)\). This factor is retained when it acts in the reader. Let

\[
C_n=e^{-x_n}M_{n+1}E_{\tau_n}(t_n)M_n^{-1}.
\]

Now \(C_n^\mathsf T\Omega C_n=e^{-2x_n}\Omega\). The same expansion gives

\begin{equation}
T_n^\mathsf T\Omega T_n+D_n^\mathsf T\Omega D_n
=\frac{8+e^{-2x_n}}9\Omega.
\label{x_reciprocidad:nuc08:eq:12}
\end{equation}

For \(x_n\geq0\), define a complementary scale register

\[
D_{{\rm esc},n}=\frac{\sqrt{1-e^{-2x_n}}}{3}B_n.
\]

The complete identity is then

\begin{equation}
\Omega=T_n^\mathsf T\Omega T_n+D_n^\mathsf T\Omega D_n
+D_{{\rm esc},n}^\mathsf T\Omega D_{{\rm esc},n}.
\label{x_reciprocidad:nuc08:eq:13}
\end{equation}

This is a composition defined here to represent explicitly the reader's scale defect. Its coefficient comes from \eqref{x_reciprocidad:nuc08:eq:12}. The name of this register does not identify it with another physical sector of the corpus; the map is retained so that this comparison can be made with their respective operators and domains.

\section{Action refinement over the complete emissions}

For each enriched prefix \(h\), let \(M_h\) be its chart and \(G_h=M_h^{-\mathsf T}M_h^{-1}\). For each surviving emission \(\epsilon\), let \(R_\epsilon\) be a realized rotation and

\[
U_\epsilon=M_{h\epsilon}R_\epsilon M_h^{-1},\qquad
(\mathscr U z)_{h\epsilon}=\sqrt{p_h(\epsilon)}U_\epsilon z_h.
\]

The weights come from the emission measure of the common construction; \(\sum_\epsilon p_h(\epsilon)=1\), with multiplicities preserved. From \(U_\epsilon^\mathsf TG_{h\epsilon}U_\epsilon=G_h\) it follows that

\begin{equation}
\sum_{h,\epsilon}I_{h\epsilon}((\mathscr U z)_{h\epsilon})
=\sum_h I_h(z_h).
\label{x_reciprocidad:nuc08:eq:14}
\end{equation}

The proof first sums the weights of a parent's descendants. Along a trajectory, the intermediate factors \(M_h^{-1}M_h\) cancel, and rotations compose in the recorded order. This preserves labels and memory without mixing children of different parents. The alternating form on the direct sum satisfies the same transport identity for the three symplectic regimes, replacing \(R_\epsilon\) with \(E_{\tau_\epsilon}(t_\epsilon)\).

The metric completions are taken with \(\sum_h z_h^\mathsf TG_hz_h\), so these maps are isometric. Further identifying this norm with a fixed Euclidean direct sum requires uniform bounds on \(M_h\) and \(M_h^{-1}\); that additional identification is not used in \eqref{x_reciprocidad:nuc08:eq:14}.
